% 乔治·伽莫夫（综述）
% license CCBYSA3
% type Wiki

本文根据 CC-BY-SA 协议转载翻译自维基百科\href{https://en.wikipedia.org/wiki/George_Gamow}{相关文章}

乔治·伽莫夫（George Gamow，有时写作 Gammoff；原名 Гео́ргий Анто́нович Га́мов，Georgiy Antonovich Gamov，1904 年 3 月 4 日－1968 年 8 月 19 日）是一位苏联和美国的博学家、理论物理学家与宇宙学家。他是乔治·勒梅特大爆炸理论的早期倡导者与发展者之一。伽莫夫首先通过量子隧穿理论解释了α衰变的机理，提出了液滴模型——第一个原子核的数学模型，并研究了放射性衰变、恒星形成、恒星核合成、以及大爆炸核合成（他统称为“核宇宙成因论”，nucleocosmogenesis），并预言了宇宙微波背景辐射的存在以及分子遗传学的基础。伽莫夫是量子隧穿现象发展与理解的关键人物之一。

在中晚期生涯中，伽莫夫将大量精力投入教学，并撰写了多部科普著作，包括《一二三……无穷大》以及“汤普金斯先生”系列丛书（1939–1967）。他的一些著作在最初出版半个多世纪后仍在印行。为纪念他，美国科罗拉多大学博尔德分校设立了乔治·伽莫夫纪念讲座。
\subsection{早年生平与职业生涯}
伽莫夫出生于俄国帝国敖德萨（今乌克兰敖德萨）。他的父亲在中学教授俄语与文学，母亲则在女子学校教授地理与历史。除了俄语外，伽莫夫还从母亲那里学会了一些法语，并由家庭教师学习了德语。在大学期间，他又学习了英语，并最终能流利使用。伽莫夫早期的大部分论文以德语或俄语撰写，后来则以英语发表技术论文和科普著作。

他先后就读于敖德萨物理与数学学院（1922–1923），以及列宁格勒大学（1923–1929）。在列宁格勒期间，伽莫夫师从亚历山大·弗里德曼，直到弗里德曼于1925年早逝，被迫更换导师。在大学时期，伽莫夫结识了三位理论物理专业的同学——列夫·朗道、德米特里·伊万年科和马特维·布龙施泰因。四人组成了一个自称为“三剑客”的小组，定期聚会讨论并分析当时发表的开创性量子力学论文。后来，伽莫夫也用同样的称呼来形容他与拉尔夫·阿尔费和罗伯特·赫尔曼共同组成的研究小组。

毕业后，伽莫夫在哥廷根从事量子理论研究，他对原子核结构的研究成果成为其博士论文的基础。此后，他于1928年至1931年在哥本哈根大学理论物理研究所工作，中途曾前往**剑桥大学卡文迪许实验室**，与著名物理学家欧内斯特·卢瑟福合作。在此期间，伽莫夫继续研究原子核（提出了著名的“**液滴模型**”），同时还与**罗伯特·阿特金森（Robert Atkinson）**和**弗里茨·豪特曼斯（Fritz Houtermans）**合作，研究**恒星物理学**。

1931年，年仅28岁的伽莫夫被选为**苏联科学院通讯院士**——成为当时最年轻的院士之一。[5][a][b] 在1931年至1933年间，他在**列宁格勒镭研究所**物理部工作，该部门由**维塔利·赫洛平（Vitaly Khlopin）**领导。在伽莫夫、**伊戈尔·库尔恰托夫与**列夫·米索夫斯基的指导与直接参与下，**欧洲第一台回旋加速器**的设计工作得以展开。1932年，伽莫夫与米索夫斯基向镭研究所学术委员会提交了加速器设计草案，并获得批准。然而，这台回旋加速器直到1937年才最终建成。[6]
\subsection{放射性衰变}
20世纪早期，科学家已经知道放射性物质具有特征性的指数衰减速率，即所谓的半衰期。与此同时，人们也发现放射性辐射具有特定的能量谱。到1928年时，伽莫夫在哥廷根成功用量子隧穿理论**解释了**原子核的α衰变现象，并在数学家尼古拉·科钦（的协助下完成了严格的推导。[7][8] 这一问题同时也被**罗纳德·W·格尼（Ronald W. Gurney）**与**爱德华·U·康登（Edward U. Condon）**独立解决。[9][10] 然而，格尼与康登的结果未能达到伽莫夫那样的定量精度。

在**经典物理学**中，粒子被束缚在原子核内，是因为要逃离强大的核势阱所需能量极高。按经典观点，必须施加巨大的外力才能将原子核分离，因此这种过程不可能自发发生。
然而，在**量子力学**框架下，粒子存在一定的**概率**可以“**穿透**”势阱的势垒并逃逸。伽莫夫在理论上求解了一个原子核的模型势场，并从**第一性原理**出发推导出**α衰变的半衰期**与**放射能量**之间的定量关系，这一关系此前是通过实验经验总结得到的，即著名的**盖革–纳塔尔定律（Geiger–Nuttall law）**。[11]

若干年后，学界将粒子穿透**库仑势垒**（Coulomb barrier）并发生核反应的概率称为**伽莫夫因子（Gamow factor）**或**伽莫夫–索末菲因子（Gamow–Sommerfeld factor）**，以纪念他在量子隧穿理论发展中的奠基性贡献。
